% Source printed pp. 117-122, PDF pp. 125-130.
\section{对角线的可表性}\label{src:4.2-section}
\begin{epigraph}
斜向运动赋予舞者一种仅有前后运动无法获得的三维感. 沿斜向运动时, 自然会有更多动作呈现在视野中.
\epigraphauthor{Lucinda Childs}
\end{epigraph}
我们证明代数空间和代数叠的对角线必然可表. 对角线用于定义稳定子和惯性叠.
\subsection{可表性}\label{src:4.2.1-subsection}
\begin{theorem}[对角线的可表性]\label{src:4.2.1}
\begin{enumerate}[(1)]
\item 代数空间的对角线可由概形表出.
\item 代数叠的对角线可表.
\end{enumerate}
\end{theorem}
\begin{proof}
设 $X$ 是代数空间, $U\to X$ 是平展呈示. 定义概形 $R:=U\times_X U$. 若 $T\to X\times X$ 是来自概形的态射, 则需证明层 $Q_T=X\times_{X\times X}T$ 实际上是概形. 由于 $U\to X$ 平展, 满射且可由概形表出, $U\times U\to X\times X$ 亦然. 将 $T\to X\times X$ 沿 $U\times U\to X\times X$ 作基变换, 得到概形 $T'$, 它在 $T$ 上满射且平展. 在 Cartesian 立方体
\begin{equation}\label{src:4.2.2}
\begin{tikzcd}[row sep=1.1em,column sep=1.4em]
{}&Q_{T'}\arrow[rr]\arrow[dl]\arrow[dd]&{}&T'\arrow[dl]\arrow[dd]\\
R\arrow[rr]\arrow[dd]&{}&U\times U\arrow[dd]&{}\\
{}&Q_T\arrow[rr]\arrow[dl]&{}&T\arrow[dl]\\
X\arrow[rr]&{}&X\times X,&{}
\end{tikzcd}
\end{equation}
中, $Q_T$ 是 $\Sch_{\text{ét}}$ 上的层, 而 $Q_{T'}$ 是概形. 由于 $R\to U\times U$ 是局部拟有限且分离的概形态射, $Q_{T'}\to T'$ 亦然. (若 $X$ 的对角线拟紧, 则由 Zariski 主定理, $R\to U\times U$ 拟仿射, 因而 $Q_{T'}\to T'$ 也拟仿射.) 由于 $Q_T$ 是平展拓扑中的层, 拉回后成为在 $T'$ 上局部拟有限且分离的概形 $Q_{T'}$, 可以应用 fppf 层成为概形的下降判据 (\bookref{3.3.17}), 得出 $Q_T$ 是概形.

若 $\mathcal X$ 是代数叠, $U\to\mathcal X$ 是光滑呈示, 则可以仿照上述论证. 纤维积 $R:=U\times_{\mathcal X}U$ 是代数空间. 若 $T\to\mathcal X\times\mathcal X$ 是来自概形的态射, 沿 $U\times U\to\mathcal X\times\mathcal X$ 作基变换得到代数空间 $T_1$, 它在 $T$ 上满射且光滑. 选取平展呈示 $T_2\to T_1$. 则 $T_2\to T$ 是概形的满射光滑态射, 在平展覆盖 $T'\to T$ 后具有截面 (命题~\bookref{A.3.4}). 复合 $T'\to T_2\to T_1\to U\times U$ 给出 $T\to\mathcal X\times\mathcal X$ 的提升. 我们得到类似于~\bookeqref{4.2.2}~的图, 但左右两个方块未必是 Cartesian 方块. 由于 $Q_T$ 等同于 $\underline{\Isom}_{\mathcal X(T)}(a,a)$, 其中 $a:T\to\mathcal X$ (见习题~\bookref{3.4.39}), $\mathcal X$ 为叠的\hyperlink{stack-axiom-1}{公理 (1)} 蕴涵 $Q_T$ 是层. 态射 $Q_{T'}\to Q_T$ 平展, 满射且可由概形表出 (因为 $T'\to T$ 如此). 选取代数空间 $Q_{T'}$ 的平展呈示 $V\to Q_{T'}$, 复合 $V\to Q_{T'}\to Q_T$ 给出平展呈示, 从而说明 $Q_T$ 是代数空间.
\end{proof}
\begin{corollary}\label{src:4.2.3}
\begin{enumerate}[(1)]
\item 若叠 $\mathcal X$ 的对角线可表 (相应地, 可由概形表出), 则每个来自概形的态射 $U\to\mathcal X$ 都可表 (相应地, 可由概形表出).
\item 从概形到代数叠 (相应地, 代数空间) 的每个态射都可表 (相应地, 可由概形表出).
\end{enumerate}
\end{corollary}
\begin{proof}
第一部分由习题~\bookref{3.4.38}~中的 Cartesian 图得到:
\[
\begin{tikzcd}
T_1\times_{\mathcal X}T_2\arrow[r]\arrow[d]\arrow[dr,phantom,"\square"]&T_1\times T_2\arrow[d]\\
\mathcal X\arrow[r]&\mathcal X\times\mathcal X.
\end{tikzcd}
\]
此图对应于从概形到代数叠的任意两个映射 $T_1\to\mathcal X$ 和 $T_2\to\mathcal X$. 第二部分直接由第一部分和对角线的可表性 (定理~\bookref{4.2.1}) 得到.
\end{proof}
\begin{exercise}\label{src:4.2.4}
\begin{enumerate}[(a)]
\item 若 $\mathcal X\to\mathcal Y$ 是代数叠的可表态射 (例如, 代数空间的态射), 证明 $\mathcal X\to\mathcal X\times_{\mathcal Y}\mathcal X$ 可由概形表出.
\item 若 $\mathcal X\to\mathcal Y$ 是代数叠的态射, 证明 $\mathcal X\to\mathcal X\times_{\mathcal Y}\mathcal X$ 可表.
\end{enumerate}
\end{exercise}
\subsection{稳定子群与惯性叠}\label{src:4.2.2-subsection}
既然知道了对角线可表, 就可以讨论它的性质. 关键在于, 对角线编码了自同构群或稳定子群. 稳定子群.
\begin{definition}[稳定子]\label{src:4.2.5}
若 $\mathcal X$ 是代数叠, $x:\Spec k\to\mathcal X$ 是域值点, 则 $x$ 的\newterm{稳定子}{stabilizer} (或 $x$ 的自同构群) 定义为群代数空间 $G_x:=\underline{\Aut}_{\mathcal X(k)}(x)$.
\end{definition}
由习题~\bookref{3.4.39}, 稳定子 $G_x$ 等同于纤维积
\[
\begin{tikzcd}
G_x=\underline{\Aut}_{\mathcal X(k)}(x)\arrow[r]\arrow[d]\arrow[dr,phantom,"\square"]&\Spec k\arrow[d,"{(x,x)}"]\\
\mathcal X\arrow[r,"\Delta"']&\mathcal X\times\mathcal X,
\end{tikzcd}
\]
由对角线的可表性 (\bookref{4.2.1}), 它可由 $k$ 上的代数空间表出. 稳定子 $G_x$ 是群代数空间, 即带有乘法, 逆元和单位态射的代数空间 $G_x$, 这些态射满足定义~\bookref{B.1.1}~中的交换条件 (等价地, 它是代数空间范畴中的群对象). 实际上, 只要 $\mathcal X$ 的对角线拟分离, $G_x$ 就是局部有限型群概形 (推论~\bookref{5.5.31}).
\begin{remark}\label{src:4.2.6}
设域 $k$ 上的群概形 $G$ 通过 $\sigma:G\times U\to U$ 作用于 $k$-概形 $U$, 且 $u\in U(k)$. $u$ 在 $[U/G]$ 中的像的稳定子是通常的稳定子群概形, 即 $(\sigma,p_2):G\times U\to U\times U$ 沿 $(u,u):\Spec k\to U\times U$ 的纤维积.
\end{remark}
\begin{exercise}\label{src:4.2.7}
\begin{enumerate}[(a)]
\item 证明代数叠的纤维积的域值点的稳定子是稳定子的纤维积, 即对于 $x'\in(\mathcal X\times_{\mathcal Y}\mathcal Y')(k)$, 有 $G_{x'}\cong G_x\times_{G_y}G_{y'}$, 其中 $x$, $y$ 和 $y'$ 是 $x'$ 的像.
\item 设 $f:\mathcal X\to\mathcal Y$ 是代数叠的态射, $x\in\mathcal X(k)$ 是域值点. 证明对角线 $\mathcal X\to\mathcal X\times_{\mathcal Y}\mathcal X$ 在点 $(x,x,\id)\in(\mathcal X\times_{\mathcal Y}\mathcal X)(k)$ 上的纤维等同于 $\ker(G_x\to G_y)$. 对角线在 $\mathcal X\times_{\mathcal Y}\mathcal X$ 的任意域值点上的纤维是什么?
\end{enumerate}
\end{exercise}
\begin{exercise}\label{src:4.2.8}
设 $\mathcal X$ 是 Deligne--Mumford 叠.
\begin{enumerate}[(a)]
\item 对于域值点 $x\in\mathcal X(k)$, 证明 $G_x$ 是 $k$ 上分离的平展群概形. 若 $G_x$ 还拟紧\footnote{这可由 $\mathcal X$ 的拟分离性推出; 见定义~\bookref{4.3.11}.}, 证明 $G_x$ 是 $k$ 上的有限平展群概形.

\emph{提示:} 先证明 $G_x$ 是 $k$ 上的平展群代数空间, 且沿某个有限域扩张 $k\to k'$ 作基变换后成为概形. 应用命题~\bookref{B.1.9}~得出 $G_x\times_k k'$ 分离. 再应用 fppf 层成为概形的下降判据 (\bookref{3.3.17}).
\item 证明 $\mathcal X$ 的对角线非分歧.\footnote{由于概形的非分歧态射在基变换下稳定, 且在源上是平展局部的性质, 此概念对代数叠的可表态射有良好定义.}
\end{enumerate}
稍后将看到, 这些性质刻画了 Deligne--Mumford 叠; 见定理~\bookref{4.6.4}.
\end{exercise}
\begin{remark}\label{src:4.2.9}
若 $\mathcal X$ 是拟分离的 Deligne--Mumford 叠, 上一习题蕴涵 $x\in\mathcal X(k)$ 的稳定子 $G_x$ 是有限平展群概形. 当 $k$ 代数闭时, $G_x$ 是抽象有限群 $\Aut_{\mathcal X(k)}(k)$ 对应的有限群概形.
\end{remark}
稳定子群 $G_x$ 随 $x\in\mathcal X$ 变化, 并自然构成一个族. 我们已经见过这一点: 若 $a:T\to\mathcal X$ 是对象, 则 $\underline{\Isom}_{\mathcal X(T)}(a)\to S$ 是群代数空间, 在点 $s\in S$ 上的纤维是 $a$ 在 $\Spec\kappa(s)$ 上的限制 $a|_{\Spec\kappa(s)}$ 的稳定子. 将此应用于恒等映射 $\id_{\mathcal X}:\mathcal X\to\mathcal X$, 就得到惯性叠.
\begin{definition}[惯性叠]\label{src:4.2.10}
代数叠 $\mathcal X$ 的\newterm{惯性叠}{inertia stack} 是纤维积
\[
\begin{tikzcd}
I_{\mathcal X}\arrow[r]\arrow[d]\arrow[dr,phantom,"\square"]&\mathcal X\arrow[d]\\
\mathcal X\arrow[r]&\mathcal X\times\mathcal X.
\end{tikzcd}
\]
在代数叠态射 $\mathcal X\to\mathcal Y$ 的相对情形中, \newterm{相对惯性叠}{relative inertia stack} 是 $I_{\mathcal X/\mathcal Y}:=\mathcal X\times_{\mathcal X\times_{\mathcal Y}\mathcal X}\mathcal X$.
\end{definition}
惯性叠 $I_{\mathcal X}$ 等价于二元组 $(x,\alpha)$ 构成的范畴, 其中 $x\in\mathcal X$ 是位于 $S\in\Sch$ 上的对象, $\alpha\in\Aut_{\mathcal X(S)}(x)$ (见习题~\bookref{4.2.15}(a)). $I_{\mathcal X}\to\mathcal X$ 在域值点 $x\in\mathcal X(k)$ 上的纤维恰好是稳定子 $G_x$. 因此可以将 $I_{\mathcal X}$ 看作 $\mathcal X$ 上囊括了 $\mathcal X$ 的所有稳定子的群概形 (更确切地说, 是群代数空间). 若 $(\Sch/\mathcal X)_{\text{ét}}$ 是 $\mathcal X$ 上概形的大平展景, 则 $I_{\mathcal X}$ 是 $(\Sch/\mathcal X)_{\text{ét}}$ 上的群层, 对于 $a\in\mathcal X(S)$, 有 $I_{\mathcal X}(a)=\Aut_S(a)$.

若 $\alpha:a'\to a$ 是位于 $S'\to S$ 上的态射, 则有自然的拉回函子 $\alpha^*:\Aut_S(a)\to\Aut_{S'}(a')$, 定义如下: 对于 $\beta\in\Aut_S(a)$, 像 $\alpha^*(\beta)$ 是使图
\begin{equation}\label{src:4.2.11}
\begin{tikzcd}[column sep=large]
a'\arrow[r,dashed,"\alpha^*(\beta)"']\arrow[rr,bend left=35,"\beta\circ\alpha"]&a'\arrow[r,"\alpha"']&a
\end{tikzcd}
\end{equation}
交换的唯一虚线箭头 (由预叠的公理 (2), 即定义~\bookref{3.4.1}, 保证). 注意, 若 $\alpha:a'\to a$ 是同构, 则 $\alpha^*(\beta)=\alpha^{-1}\circ\beta\circ\alpha$ 是由 $\alpha$ 给出的共轭.
\begin{exercise}\label{src:4.2.12}
设 $G\to S$ 是作用于概形 $U\to S$ 的群概形, $\mathcal X=[U/G]$ 是商叠. 证明存在 Cartesian 图
\[
\begin{tikzcd}
S_U\arrow[r]\arrow[d]\arrow[dr,phantom,"\square"]&U\arrow[d]\\
I_{\mathcal X}\arrow[r]&\mathcal X,
\end{tikzcd}
\]
其中 $S_U\to U$ 是稳定子群概形, 即作用映射 $G\times U\to U\times U$ 与对角线 $U\to U\times U$ 的纤维积.
\end{exercise}
\begin{example}\label{src:4.2.13}
分类叠 $B\mathbb G_m$ 的惯性叠是 $I_{B\mathbb G_m}\cong\mathbb G_m\times B\mathbb G_m$. 类似地, 令 $\mathbb G_m$ 通过平凡作用与数乘作用的乘积作用于 $\mathbb G_m\times\mathbb A^1$, 并令 $V(x(t-1))\subseteq\mathbb G_m\times\mathbb A^1$ 为 $\mathbb G_m$-不变闭子概形, 则 $I_{[\mathbb A^1/\mathbb G_m]}\cong[V(x(t-1))/\mathbb G_m]$.
\end{example}
\begin{exercise}\label{src:4.2.14}
\begin{enumerate}[(a)]
\item 若 $G$ 是域 $k$ 上的光滑仿射代数群, 证明 $BG$ 的惯性叠是商 $[G/G]$, 其中 $G$ 通过共轭作用于自身. 特别地, 若 $G$ 交换, 则 $I_{BG}\cong G\times BG$.
\item 更一般地, 若 $G$ 作用于 $k$-概形 $U$, 证明 $I_{[U/G]}\cong[(G\times U)/G]$, 其中作用由 $g\cdot(h,u)=(ghg^{-1},gu)$ 给出.
\item 设 $G$ 是 $k$ 上对应于某个有限抽象群的群概形. 若 $G$ 作用于 $k$-概形 $U$, 则
\[
I_{[U/G]}=\coprod_{g\in\operatorname{Conj}(G)}[U^g/C_g],
\]
其中 $\operatorname{Conj}(G)$ 是 $G$ 中元素的共轭类集合, $C_g$ 是 $g$ 的中心化子, $U^g:=\{x\in U\mid gx=x\}$ 是被 $g$ 固定的闭子概形.
\item 显式计算置换作用的商叠 $[\mathbb A^3/S_3]$ 的惯性叠.
\end{enumerate}
\end{exercise}
% Keep the exercise heading with its first item, rather than alone at the foot.
\newpage
\begin{exercise}\label{src:4.2.15}
设 $f:\mathcal X\to\mathcal Y$ 是代数叠的态射.
\begin{enumerate}[(a)]
\item 证明相对惯性叠 $I_{\mathcal X/\mathcal Y}$ 等价于二元组 $(x,\alpha)$ 构成的范畴, 其中 $x\in\mathcal X$, $\alpha:x\xrightarrow{\sim}x$ 是满足 $f(\alpha)=\id_{f(x)}$ 的同构; 并证明单位截面 $\mathcal X\to I_{\mathcal X/\mathcal Y}$ 将对象 $x$ 送到 $(x,\id_x)$.
\item 证明存在 $\mathcal X$ 上的代数叠态射 $I_{\mathcal X/\mathcal Y}\to I_{\mathcal X}\to I_{\mathcal Y}\times_{\mathcal Y}\mathcal X$, 它们在域值点 $x\in\mathcal X(k)$ 上对应于群代数空间的左正合列 $1\to K_x\to G_x\to G_{f(x)}$.
\item 证明存在 Cartesian 图
\[
\begin{tikzcd}
I_{\mathcal X}\arrow[r]\arrow[d]\arrow[dr,phantom,"\square"]&I_{\mathcal Y}\times_{\mathcal Y}\mathcal X\arrow[d]\\
\mathcal X\arrow[r]&\mathcal X\times_{\mathcal Y}\mathcal X.
\end{tikzcd}
\]
\end{enumerate}
\emph{(c) 的提示:} $I_{\mathcal Y}\times_{\mathcal Y}\mathcal X$ 中位于概形 $S$ 上的对象是四元组 $(y,\alpha,x,\beta)$, 其中 $y\in\mathcal Y(S)$, $\alpha:y\xrightarrow{\sim}y$, $x\in\mathcal X(S)$, $\beta:y\xrightarrow{\sim}f(x)$. 另一方面, $\mathcal X\times_{\mathcal Y}\mathcal X$ 中位于 $S$ 上的对象是三元组 $(x_1,x_2,\gamma)$, 其中 $x_1,x_2\in\mathcal X(S)$, $\gamma:f(x_1)\xrightarrow{\sim}f(x_2)$. 在纤维范畴上通过 $(y,\alpha,x,\beta)\mapsto(x,x,\beta\circ\alpha\circ\beta^{-1})$ 定义 $I_{\mathcal Y}\times_{\mathcal Y}\mathcal X\to\mathcal X\times_{\mathcal Y}\mathcal X$. 构造从 $I_{\mathcal X}(S)$ 到 $\mathcal X(S)$ 与 $(I_{\mathcal Y}\times_{\mathcal Y}\mathcal X)(S)$ 在 $(\mathcal X\times_{\mathcal Y}\mathcal X)(S)$ 上的纤维积的映射, 并证明它是等价.
\end{exercise}
